\section{CLIP: Contrastive Language–Image Pretraining}
\medskip

So far, we have learned how to embed text into a vector space. We now ask a broader question: can we build a single embedding space that contains both text and images?

Suppose we want to search an image database using natural language. A user types the query ``panda.'' We would like the system to return images of pandas. One naive solution would be to train an image classifier with a fixed set of categories, for example the 1,000 ImageNet labels, classify every image, and then look for the label ``panda.'' This works only if the category exists in the label set. If the user searches for something more specific or unusual, such as ``sleepy panda eating bamboo in snow,'' the fixed label taxonomy breaks down.

The core limitation is that classification assumes a closed set of categories. Real-world search is open-ended. We need a system that understands arbitrary text queries and arbitrary images.

\bigskip
\subsection{Where Does the Data Come From?}
\medskip

Training large image classifiers traditionally requires hand-labeled data. ImageNet contains roughly 15 million images with carefully curated labels. While extremely valuable, this approach does not scale indefinitely. Labeling is expensive, slow, and constrained to predefined categories.

However, the internet already contains an enormous number of image–text pairs. When people post on social media, they often include a caption alongside an image. The caption may be formal, casual, noisy, or emotional, but it is text paired with an image. In aggregate, this provides supervision at scale.

The key idea behind CLIP is to treat these naturally occurring image–text pairs as training data.

\bigskip
\subsection{Two Encoders, One Space}
\medskip

The model architecture mirrors what we have already seen in bi-encoders for retrieval.

We build:
\begin{itemize}
    \item A text encoder that maps a sentence to a vector.
    \item An image encoder that maps an image to a vector.
\end{itemize}

Given $N$ paired examples $(T_i, I_i)$, we compute embeddings:
\[
t_i = f_\theta(T_i), 
\qquad
v_i = g_\phi(I_i).
\]

The goal is to learn parameters $\theta$ and $\phi$ such that matching text–image pairs are close in vector space, and mismatched pairs are far apart.

Intuitively, we want both modalities to ``speak the same geometric language.'' A picture of a panda and the word ``panda'' should land near each other in $\mathbb{R}^d$.

\bigskip
\subsection{The Contrastive Objective}
\medskip

Training proceeds using a contrastive loss, very similar to the InfoNCE loss we discussed for text retrieval.

For a batch of size $N$, we compute all pairwise similarities:
\[
S_{ij} = \langle t_i, v_j \rangle.
\]

Each text $T_i$ should match only its paired image $I_i$. For the text-to-image direction, the loss is:
\[
\mathcal{L}_i^{\text{text}}
=
- \log
\frac{\exp(S_{ii})}
{\sum_{j=1}^{N} \exp(S_{ij})}.
\]

Symmetrically, we also treat images as queries and texts as candidates:
\[
\mathcal{L}_i^{\text{image}}
=
- \log
\frac{\exp(S_{ii})}
{\sum_{j=1}^{N} \exp(S_{ji})}.
\]

The final loss averages both directions.

What is happening geometrically? For each pair $(T_i, I_i)$, we pull $t_i$ and $v_i$ together, while pushing them away from all other embeddings in the batch. The batch itself provides negative examples. No manual labels are required beyond the pairing.

\bigskip
\subsection{Geometric Intuition}
\medskip

It is helpful to visualize embeddings as points on a high-dimensional sphere. During training, matching text and image vectors are encouraged to align, meaning their cosine similarity is high. Non-matching pairs are encouraged to point in different directions.

The loss is small when:
\begin{itemize}
    \item $t_i$ is close to $v_i$,
    \item $t_i$ is far from $v_j$ for $j \neq i$.
\end{itemize}

This is exactly the same contrastive principle we saw in bi-encoder retrieval, now extended across modalities.

\bigskip
\subsection{Why This Is Powerful}
\medskip

After training, text and images share a single embedding space. This enables several capabilities:

\begin{itemize}
    \item Text-to-image search: embed a text query and retrieve nearest images.
    \item Image-to-text search: embed an image and retrieve nearest captions.
    \item Zero-shot classification: treat class names as text queries and select the closest label embedding.
\end{itemize}

Zero-shot classification is particularly striking. Instead of training a classifier for each dataset, we can embed candidate labels such as ``panda,'' ``tiger,'' or ``hedgehog'' and choose the one whose text embedding is closest to the image embedding. The model generalizes to categories it has never explicitly seen during supervised training.

\bigskip
\subsection{Connection to Earlier Ideas}
\medskip

Conceptually, CLIP combines ideas we have already encountered:

\begin{itemize}
    \item A bi-encoder architecture.
    \item Contrastive learning via an InfoNCE-style loss.
    \item Large-scale self-supervision using naturally occurring data.
\end{itemize}

The only difference is that the two encoders operate on different modalities. The objective remains the same: align related representations and separate unrelated ones.

In this sense, CLIP extends embedding learning from text-only retrieval to multimodal retrieval. The geometry of the space becomes a shared semantic space, where images and language coexist.

The central insight is simple but powerful: if two objects co-occur in the world, their embeddings should be close. Contrastive learning gives us the mathematical machinery to enforce this principle at scale.


\bigskip

